\subsection{A spectral resolution on the line}
\label{subsec:bethe-line}

The ground-state identification will start at strong coupling, where each
singlet pair has binding energy \(-c^2/2\).  We need to know that every
state with energy \(-Mc^2/2+O(1)\) is concentrated on \(M\) such pairs,
with at most one single particle.  We prove this first on the line by a
spectral resolution, and then pass to a fixed circle by localization.
Throughout this argument the particle number \(n_p=2M+\ell\) is fixed,
with \(\ell\in\{0,1\}\), and the coupling \(c>0\) is not scaled to \(2\).

Write
\[
 \mathcal C_{n_p}=\{X\in\mathbb R^{n_p}:X_1<\cdots<X_{n_p}\},
 \qquad
 \mathcal V_{n_p}=(\mathbb C^2\otimes\mathbb C^2)^{\otimes n_p}.
\]
We realize the line Hamiltonian \(H_c\) on
\(L^2(\mathcal C_{n_p};\mathcal V_{n_p})\), with the contact conditions
already derived from \eqref{eq:B2}.  Lebesgue measure on the ordered
chamber is used throughout.  For two slots \(a,b\), let
\(s_{ab}=2^{-1/2}(e_1\otimes e_2-e_2\otimes e_1)\) be the normalized
singlet in one binary index, and let
\(s_{ab}^{L}\otimes s_{ab}^{R}\) be the double singlet.

For integers \(\ell',m'\geq0\) satisfying \(\ell'+2m'=n_p\), choose
\(\ell'\) real elementary momenta \(k_1,\ldots,k_{\ell'}\) and \(m'\)
real pair centers \(u_1,\ldots,u_{m'}\).  The two momenta belonging to
center \(u_a\) are
\[
                 u_a-\frac{ic}{2},\qquad u_a+\frac{ic}{2}.
\]
Initially put each pair in adjacent slots, with its lower momentum
first, and insert its double singlet.  The elementary slots retain
arbitrary colors.  The resulting plane-wave sum, denoted
\(E_X^{\ell'm'}(k,u)\), is an operator
\(\mathcal V_{\ell'}\to\mathcal V_{n_p}\).  Its coefficients are obtained
with the interchange matrix \(\check R\); the sum runs over all
shuffles preserving the order of the two members of every pair.
More explicitly, let \(\kappa_1,\ldots,\kappa_{n_p}\) be this initial
list of momenta, and let \(\iota:\mathcal V_{\ell'}\to\mathcal V_{n_p}\)
insert the normalized double singlets in its pair slots.
If \(\mathfrak S_{\ell'm'}\) consists of the permutations \(\pi\)
of labels into slots satisfying
\(\pi(A_a)<\pi(B_a)\) for every labeled low--high pair \(A_aB_a\),
then
\[
 E_X^{\ell'm'}(k,u)
   =\sum_{\pi\in\mathfrak S_{\ell'm'}}
       e^{\,i\sum_{j=1}^{n_p}\kappa_jX_{\pi(j)}}A_\pi(k,u)\iota,
\]
where \(A_\pi\) is the product of adjacent interchange matrices
along a reduced sequence producing \(\pi\), and \(A_{\rm id}=I\).
Coefficients reversing a pair are zero.  The inverse and braid
relations make this prescription independent of the chosen sequence
of adjacent interchanges.  At a removable singularity we use its
continuous value.  The initial order of the whole blocks can be chosen
arbitrarily; the block-interchange formulas proved below describe the
change of this convention.

\begin{lemma}[Line spectral resolution]
\label{lem:bethe-line}
For every fixed integer \(n_p=2M+\ell\geq0\), every \(c>0\), and every
\(s>0\), the heat kernel of \(H_c\) on the ordered line chamber is
\begin{equation}
\begin{split}
 e^{-sH_c}(X,Y)
  ={}&\sum_{\ell'+2m'=n_p}
       \frac{(2c)^{m'}}{\ell'!\,m'!}
       \int_{\mathbb R^{\ell'+m'}}
       E_X^{\ell'm'}(k,u)E_Y^{\ell'm'}(k,u)^\dagger\\
    &\quad{}\times
       \exp\left\{-s\left(
           \sum_{a=1}^{\ell'}k_a^2+
           2\sum_{a=1}^{m'}u_a^2-\frac{m'c^2}{2}
                    \right)\right\}
       \prod_{a=1}^{\ell'}\frac{dk_a}{2\pi}
       \prod_{a=1}^{m'}\frac{du_a}{2\pi}.
\end{split}
\label{eq:B10}
\end{equation}
Empty products and the vacuum term have their usual value \(1\).
The transform defined by the kernels in \eqref{eq:B10}, with the
square roots of the displayed weights, is an isometry.  Moreover,
\(H_c+Mc^2/2\geq0\), and its quadratic form is the integral of
\[
  \sum_{a=1}^{\ell'}k_a^2+
  2\sum_{a=1}^{m'}u_a^2+\frac{(M-m')c^2}{2}
\]
against the squared modulus of this transform.
\end{lemma}

\begin{proof}
The proof has three parts.  The scattering identities show that taking
a pair residue creates no further bound cluster.  A heat-kernel formula
on separated contours then supplies the initial condition.  Finally,
lowering those contours produces \eqref{eq:B10}, with its normalization
determined by the particle occupying the first slot at each step.

\emph{Scattering of pairs and cancellation of higher residues.}
Let \(\Pi_{\varepsilon\eta}\), with
\(\varepsilon,\eta\in\{+1,-1\}\), be the joint spectral projections of
\(P^L,P^R\) on two slots.  The interchange matrix has the decomposition
\[
 \check R(x)=
 \frac{x-ic}{x+ic}\Pi_{++}
 +\frac{x+ic}{x-ic}\Pi_{--}
 -\Pi_{+-}-\Pi_{-+}.
\]
In particular,
\[
 \mathop{\rm Res}_{x=ic}\check R(x)=2ic\,\Pi_{--},
 \qquad \check R(0)=-I.
\]
Here \(\Pi_{--}\) is precisely the double-singlet projection.  Taking
this residue produces a low--high pair in the order specified above.
Multiplication of two interchanges transports a pair past a single
and multiplies its color vector by
\begin{equation}
 S_{21}(x)=
 \frac{(x+ic/2)(x-3ic/2)}
      {(x-ic/2)(x+3ic/2)}.
\label{eq:bethe-pair-single}
\end{equation}
In this formula \(x\) is the center of the left block minus the
momentum of the right block.  If the single is on the left, the same
formula holds with \(x\) equal to its momentum minus the pair center.
Moving a pair past another pair gives
\begin{equation}
                       S_{22}(x)=\frac{x-2ic}{x+2ic},
\label{eq:bethe-pair-pair}
\end{equation}
where again \(x\) is the left center minus the right center.
Both scalars have modulus one for real \(x\).

Here are explicit checks, including the partial exchanges that will
occur during contour motion.  In one binary index, for a pair \(AB\)
followed by a single \(C\), set
\[
 E(q)=s_{12}\otimes q_3,\qquad
 H(q)=q_1\otimes s_{23},\qquad
 F(q)=s_{13}\otimes q_2=E(q)+H(q).
\]
For an operator or vector \(V\) in one binary index, use
\(V^{[2]}=V^L\otimes V^R\).  Put
\[
 D(x)=(x-ic/2)(x+3ic/2),\qquad
 T(x)=(x-ic/2)E+(x+ic/2)H.
\]
With \(x=u-k\), the coefficients of the three shuffles \(ABC,ACB,CAB\)
are, respectively,
\[
            E^{[2]},\qquad
            \frac{T(x)^{[2]}}{D(x)},\qquad
            S_{21}(x)H^{[2]}.
\]
These identities follow by applying the two permutation factors in
\(\check R\) to the singlet.  For a single initially on the left,
interchange \(E\) and \(H\), and take \(x=k-u\).  Thus neither partial
nor complete exchanges have a pole at \(x=3ic/2\).  At \(x=ic/2\),
the last two residues in the pair-first order are
\((ic/2)H^{[2]}\) and \(-(ic/2)H^{[2]}\).
The elementary momentum then equals the lower member of the pair,
so their plane phases agree and the residues cancel.  In the
single-first order the corresponding equality is with the upper
member.  This is also the cancellation expressed by
\(\check R(0)=-I\).

For two pairs \(AB,CD\), both in low--high order, put \(x=u-v\).
In one binary index set
\[
        e=s_{12}s_{34},\qquad f=s_{14}s_{23},\qquad
        h=s_{13}s_{24}=e+f.
\]
The coefficients of the six allowed shuffles are
\[
\begin{array}{c|c}
 ABCD & e^{[2]}\\[2pt]
 ACBD & (xh+icf)^{[2]}/\{x(x+2ic)\}\\[2pt]
 ACDB & (x^2+c^2)f^{[2]}/\{x(x+2ic)\}\\[2pt]
 CABD & (x^2+c^2)f^{[2]}/\{x(x+2ic)\}\\[2pt]
 CADB & (xh-icf)^{[2]}/\{x(x+2ic)\}\\[2pt]
 CDAB & (x-2ic)e^{[2]}/(x+2ic).
\end{array}
\]
This table can equally be obtained from \(P_{12}=-1\) on a singlet
and
\(1-P_{12}-P_{13}-P_{23}+P_{12}P_{23}+P_{23}P_{12}=0\)
on three binary slots.  It proves \eqref{eq:bethe-pair-pair} and
shows that the apparent poles at \(x=ic\) and \(x=2ic\) are absent.
At \(x=0\), the four middle residues are, in their listed order,
\[
       \frac{ic}{2}f^{[2]},\quad
       -\frac{ic}{2}f^{[2]},\quad
       -\frac{ic}{2}f^{[2]},\quad
       \frac{ic}{2}f^{[2]}.
\]
The phases agree in the pairs \(ACBD,ACDB\) and \(CABD,CADB\),
because \(B\) and \(D\) have the same momentum.  The cancellation
therefore holds both in the whole wave and in either part in which
the first slot is required to contain the low member of a specified
pair: the part starting with \(A\) consists of the first three
shuffles, and the part starting with \(C\) of the last three.

\emph{The heat kernel on separated contours.}
Let \(E_X^{n_p,0}(k)\) be the elementary-wave matrix, and choose
horizontal contours
\(\Gamma_j=\mathbb R+i\alpha_j\), where
\[
 \alpha_{n_p}=0,\qquad \alpha_j-\alpha_{j+1}>c.
\]
Consider
\begin{equation}
 \mathcal K_s(X,Y)=
 \int_{\Gamma_1\times\cdots\times\Gamma_{n_p}}
 E_X^{n_p,0}(k)
 e^{-i\sum_j k_jY_j-s\sum_j k_j^2}
 \prod_j\frac{dk_j}{2\pi}.
\label{eq:bethe-separated-contours}
\end{equation}
Gaussian decay makes this integral and its spatial derivatives
convergent on compact sets.  Each plane wave solves the free heat
equation in the chamber, and the interchange relation imposes the
contact conditions on each face.

To verify the initial condition, one needs the support of the
Fourier transforms of the rational coefficients.  For
\(\operatorname{Im}x>c\), the preceding projector decomposition gives
the concrete formula
\[
 \check R(x)=P^LP^R+
 2c\int_0^\infty e^{ixt}
       \bigl(e^{ct}\Pi_{--}-e^{-ct}\Pi_{++}\bigr)\,dt.
\]
Accordingly, an inversion of labels \(i<j\) contributes a
nonnegative variable \(t_{ij}\), possibly an atom at zero, and a phase
\(e^{i(k_i-k_j)t_{ij}}\).  If a permutation puts label \(i\) at
position \(p(i)\), its zero-time contribution can be supported only at
\begin{equation}
 Y_i=X_{p(i)}+
       \sum_{\substack{j>i\\p(j)<p(i)}}t_{ij}
       -\sum_{\substack{h<i\\p(h)>p(i)}}t_{hi}.
\label{eq:bethe-inversion-support}
\end{equation}
For a nonidentity permutation, let \(r\) be the first nonfixed
position, and let \(j>r\) be the label occupying it.
Label \(r\) occurs at some position \(p(r)>r\).
No smaller label contributes a negative shift to label \(r\);
no larger label contributes a positive shift to label \(j\).
Thus \eqref{eq:bethe-inversion-support} would imply
\[
                   Y_r\geq X_{p(r)}>X_r\geq Y_j,
\]
which is incompatible with \(Y\in\mathcal C_{n_p}\).
Only the identity permutation remains, and its coefficient is the
identity color operator.  It supplies the delta initial condition.

The support argument also supplies the estimates needed to justify
this limiting calculation and heat-kernel uniqueness.  Let \(d\)
be the shift vector on the right side of
\eqref{eq:bethe-inversion-support}, so that
\(z_i=X_{p(i)}+d_i\).  The shifts belong to the cone generated by
\(e_i-e_j\), \(i<j\).  This cone is pointed, and more explicitly
\[
 \sum_{i<j}t_{ij}
 \leq \sum_{i<j}(j-i)t_{ij}
 =-\sum_i i\,d_i
 \leq C_{n_p}|d|.
\]
Furthermore,
\[
 \sum_i X_{p(i)}d_i
   =\sum_{\substack{i<j\\p(j)<p(i)}}
        t_{ij}(X_{p(i)}-X_{p(j)})\geq0,
 \qquad |z|^2\geq |X|^2+|d|^2.
\]
The measures in the transform formula have at most exponential
growth in \(\sum t_{ij}\).  Gaussian convolution therefore dominates
them and gives spatial decay for compactly supported initial data.
It also justifies convergence to the stated initial condition,
first for data supported strictly inside the chamber and then by
density.  The contact form is lower bounded by the hyperplane trace
inequality, so uniqueness of its \(L^2\) heat evolution identifies
\(\mathcal K_s\) with \(e^{-sH_c}\).

\emph{Lowering the contours and reconstructing the adjoint wave.}
Process the momenta from last to first.  The induction assertion can be
stated explicitly.  After the last \(n_p-h\) slots have been processed,
sum over \(\ell'+2m'=n_p-h\), with weight
\((2c)^{m'}/(\ell'!m'!)\).  The first \(h\) elementary momenta
\(\xi_1,\ldots,\xi_h\) remain on their original high contours; the
suffix momenta \(k\) and pair centers \(u\) are real.  The wave factor
in the integrand is
\[
 E_X^{h+\ell',m'}(\xi,k,u)
 \left[I_{\mathcal V_h}\otimes
       E_{(Y_{h+1},\ldots,Y_{n_p})}^{\ell'm'}(k,u)^\dagger\right]
 e^{-i\sum_{a=1}^h\xi_aY_a}.
\]
Its Gaussian energy is
\(\sum\xi_a^2+\sum k_a^2+2\sum u_a^2-m'c^2/2\).
Here the \(X\)-wave uses the first \(h\) elementary slots followed by
the suffix blocks as its initial order; its finite meromorphic formula
also defines it for the complex variables \(\xi\).  Thus \(h=n_p\)
is \eqref{eq:bethe-separated-contours}, whereas \(h=0\) is precisely
\eqref{eq:B10}.

To pass from \(h\) to \(h-1\), decompose the new suffix wave according
to its first occupied slot.  For each elementary carrier or labeled pair
\(a\), let \(E_Y^{[a]}\) be the partial sum in which that slot contains,
respectively, the elementary momentum or the pair's low member.  These
possibilities are disjoint and exhaust the allowed shuffles, so
\(E_Y=\sum_a E_Y^{[a]}\).  The ordinary contour integral will supply
the elementary-first partial adjoints, and its residues will supply the
pair-first partial adjoints.
When block parameters \(z\) leave the real contours, an adjoint wave is
continued as the meromorphic function \(E_Y(\overline z)^\dagger\),
which agrees with \(E_Y(z)^\dagger\) for real \(z\).  The same convention
applies to each partial adjoint.

Let \(k=\xi_h\) be the next momentum.  The integral left after lowering
its contour to the real line provides the adjoint-wave terms whose
first suffix slot is an elementary momentum.  To see this with a
fixed elementary carrier, put that carrier first among the suffix
blocks.  The unitary elementary and whole-block interchanges move
the corresponding coefficient from the \(X\)-wave to the
\(Y\)-wave.  Symmetrizing the elementary variables then supplies
all elementary choices for the first slot.

The residues occur at \(k=p+ic\), where \(p\) is an elementary
momentum in the old suffix.  Put \(p\) first in that suffix before
computing the residue.  A downward contour shift contributes
\(-i\) times the residue, so the residue \(2ic\,\Pi_{--}\) contributes
the factor \(2c\) and inserts a normalized double singlet.
The momenta at \(X\) are \(p,k\), in low--high order, with center
\(u=p+ic/2\).  Lower \(u\) from this contour to the real line.
On the \(Y\) side the new first slot uses the adjointed phase of
the low member, namely the phase with the high momentum \(k\).
The other member is the old suffix variable \(p\).  The old
suffix exchanges are continued with this \(p\), and hence give
exactly the part of \(E_Y^\dagger\) with that pair's low member
first.

For completeness, the pole checks justifying both contour motions
are as follows.  When \(k\) descends past an existing pair, the
three-slot formulas have no pole at center difference \(3ic/2\);
the possible pole at \(ic/2\) cancels between the partial and
complete exchanges whose phases agree.  There is therefore no
three-particle residue.  When the newly formed center \(u\) is
lowered, the same cancellation applies to the \(X\)-wave.  In the
restricted adjoint \(Y\)-expression the low member of the new pair
has already been put first.  The crossing that remains uses only
\(p=u-ic/2\).  Its imaginary part moves from \(0\) to \(-c/2\),
and no elementary interchange pole at difference \(\pm ic\) is
crossed.  This includes an elementary momentum on the old real
contour, by continuation with an indentation if necessary.
Explicitly, the adjoint of the pair-low-first three-slot sum has
denominator \((x+ic/2)(x-3ic/2)\), with \(x=u-k_{\rm single}\),
which has no zero in \(0\leq\operatorname{Im}x\leq c/2\).
For two pairs, the four-slot table proves regularity at coincident
centers even in the restricted first-slot expression.
The only other possible center pole is at \(-2ic\) for the
\(X\)-coefficients and at \(2ic\) for the continued adjoint
coefficients, both outside the strips being traversed.

These checks remain valid with other variables inserted.  At a
generic point of a candidate pole hyperplane, use the braid and
whole-block identities to bring the two relevant blocks together.
For fixed occupied slots, inserting the other momenta multiplies
the two cancelling residues by the same nonsingular interchange
operators.  If the first slot is restricted, put its block first
before this operation.  The earlier, unprocessed contours stay
separated from all these motions.  Finally, write the finite sums
over a common rational denominator.  Removal of the poles at
generic points of every denominator hyperplane removes the
corresponding factors, so the removals extend across their
intersections.  Distinct real parts or indented contours may
therefore be used for the residue computations and subsequently
removed.  On bounded \(X,Y\), small contour circles about removable
singularities give uniform bounds for the finite wave sums.
Together with the Gaussian, these bounds justify the integrations
and deformations just made.

It remains to give every partial adjoint the same weight.  Changing
the initial block order multiplies both waves by the same unitary
block-interchange operator, so the product
\(E_X(E_Y^{[a]})^\dagger\) is independent of that choice.  We may
therefore put each first carrier in front before symmetrizing the
labels.  For a final sector \((\ell',m')\), the
elementary-first contribution comes from
\((\ell'-1,m')\) in the old suffix and has weight
\[
 \frac{(2c)^{m'}}{(\ell'-1)!\,m'!}
       =\ell'\frac{(2c)^{m'}}{\ell'!\,m'!}.
\]
It accounts for the \(\ell'\) choices of elementary first carrier.
The pair-first contribution comes from
\((\ell'+1,m'-1)\).  There are \(\ell'+1\) elementary choices
for \(p\), and the residue supplies \(2c\); its weight is therefore
\[
 2c(\ell'+1)
 \frac{(2c)^{m'-1}}{(\ell'+1)!\,(m'-1)!}
       =m'\frac{(2c)^{m'}}{\ell'!\,m'!}.
\]
It accounts for the \(m'\) choices of pair first carrier.
After symmetrization, each of the \(\ell'+m'\) partial adjoints has
weight \((2c)^{m'}/(\ell'!m'!)\).  Their sum is the full adjoint wave.
This closes the induction and proves \eqref{eq:B10}.

Define the transform explicitly by
\[
 (\mathcal F_c f)_{\ell'm'}(k,u)=
 \left(\frac{(2c)^{m'}}{\ell'!\,m'!}\right)^{1/2}
 \int_{\mathcal C_{n_p}}E_Y^{\ell'm'}(k,u)^\dagger f(Y)\,dY,
\]
initially for compactly supported \(f\).  Its target is the direct
sum of the \(L^2\) spaces with the product measures \(dk\,du/(2\pi)\)
in \eqref{eq:B10} and color space \(\mathcal V_{\ell'}\).
Writing
\[
 \lambda_{\ell'm'}(k,u)=
 \sum k_a^2+2\sum u_a^2+\frac{(M-m')c^2}{2}\geq0,
\]
the quadratic identity from \eqref{eq:B10} is
\[
 \langle f,e^{-s(H_c+Mc^2/2)}f\rangle
 =\sum_{\ell'+2m'=n_p}\int
       e^{-s\lambda_{\ell'm'}}|(\mathcal F_cf)_{\ell'm'}|^2.
\]
As \(s\downarrow0\), strong continuity on the left and monotone
convergence on the right give \(\|\mathcal F_cf\|=\|f\|\).
The transform thus extends isometrically to the full Hilbert
space.  The same identity then implies
\(\|e^{-s(H_c+Mc^2/2)}\|\leq1\), and hence
\(H_c+Mc^2/2\geq0\).  Taking the increasing limit of
\((1-e^{-s\lambda})/s\) gives, by the spectral theorem,
the stated quadratic-form identity and its domain.
The reconstruction \(\mathcal F_c^*\mathcal F_c=I\) is all that
will be needed; surjectivity onto the direct sum is unnecessary.
\end{proof}

\subsection{Pair concentration at strong coupling}
\label{subsec:bethe-strong-coupling}

The spectral resolution contains only single particles and singlet pairs;
the residue cancellations have excluded additional binding channels.
Its form identity gives a penalty \(c^2/2\) for every missing pair.
To use this information on a circle, we also need control of the spatial
shape of the pair waves, uniformly in their real center momenta.

For a shuffle contributing to \(E_X^{\ell'm'}\), draw the closed
interval between the two slots occupied by each pair.  Combine
overlapping intervals into connected components, and include every
single whose slot lies inside such a component.  A single outside
all pair intervals forms a component of one slot.  Group the
shuffles for which every component has the same set of slots and
the same set of carriers (elementary momenta and labeled pairs).
Let \(E_{X,\mathfrak g}^{\ell'm'}\) denote the sum for one such group.

\begin{lemma}[Decay across pair spans]
\label{lem:bethe-shuffle-decay}
At fixed \(n_p\), there are constants \(C_{n_p}<\infty\) and
\(a_{n_p}>0\), independent of \(c>0\) and the real momenta, such that
every group just defined satisfies
\begin{equation}
 \bigl\|E_{X,\mathfrak g}^{\ell'm'}(k,u)\bigr\|
 \leq C_{n_p}
 \exp\left\{-a_{n_p}c
       \sum_{\text{components }\mathcal I}
       \bigl(X_{\max\mathcal I}-X_{\min\mathcal I}\bigr)\right\}.
\label{eq:B11}
\end{equation}
The norm is the operator norm after summing the shuffles in the
group.
\end{lemma}

\begin{proof}
Divide all momenta by \(c\).  A shuffle's plane phase has modulus
\[
 \exp\left\{-\frac c2
       \sum_{\text{pairs }a}
           (X_{\text{high slot of }a}
                 -X_{\text{low slot of }a})\right\}.
\]
The sum of the pair spans is at least the sum of the component
spans.  Rational interchange coefficients are uniformly bounded
away from their pole hyperplanes.  Near real centers, their only
apparent singularities are coincident pair centers.  The
four-slot cancellations in the preceding proof are internal to
the present groups: switching the equal high members preserves
the slots and carriers of each overlapping component.
More explicitly, the cancelling shuffles are
$ACBD,ACDB$ and $CABD,CADB$.  Within either pair of shuffles,
the union of the spans belonging to $AB$ and $CD$ is the same
interval, from their earliest low member to their latest high member.
Inserting further slots does not change that equality.  The connected
components obtained after adjoining all other pair spans, and their
sets of carriers and included singles, are therefore unchanged.
Thus every group sum is holomorphic through these real
hyperplanes.

Here is a uniform Cauchy estimate for this holomorphic sum.
Choose a sufficiently small number depending only on \(n_p\).
Among finitely many successively much smaller scales, one can
choose a scale with no distance between real divided momenta in
the intervening annulus: there are only finitely many pairwise
distances.  Cluster the arguments at the smaller scale.
Arguments in different clusters are then separated by much more
than the chosen scale.  Within each cluster, put the variables
on small circles about their original real values,
using distinct radii separated by fixed multiples of that scale.
All radii can be kept below the initially chosen small number.
On the product of these circles, equal-center denominators are
bounded away from zero by a positive \(n_p\)-dependent constant;
denominators with nonzero imaginary offsets are also separated
from zero.  Every rational coefficient is consequently bounded
there by a constant depending only on \(n_p\).  Cauchy's formula
recovers its holomorphic group sum at the original real arguments.

To control the plane phases on those circles, factor from each
component the common phase evaluated at its leftmost position.
The carrier set of that component is fixed throughout the group,
so the factored phase is the same in every term and has modulus
one at real momenta.  Perturbing the remaining divided momenta by
the small circle radii increases their modulus by at most
\(\exp\{\varepsilon C_{n_p}c\sum_{\mathcal I}
\operatorname{span}(\mathcal I)\}\).
Choose the initial upper bound on the radii small enough that
this uses less than half the original pair-gap decay.
The resulting Cauchy estimate proves \eqref{eq:B11}.
\end{proof}

We next specify the compactness statement on a circle.  Fix its
length \(W_0>0\).  Choose finitely many seams and smooth nonnegative,
permutation-invariant cutoffs \(\vartheta_\alpha\), with
\(\sum_\alpha\vartheta_\alpha^2=1\), such that all positions in the
support of \(\vartheta_\alpha\) stay a positive distance from the
corresponding seam.  Such a choice is possible because, at fixed
particle number, among sufficiently many candidate seams at least
one stays a fixed distance from every position.  Cutting at that
seam identifies each localized function with a compactly supported
function on a line chamber.

In each chamber let \(\mathfrak p\) be a partition of its ordered
slots into \(M\) adjacent pairs and, when \(\ell=1\), one single.
For a pair occupying slots \(a,a+1\), use coordinates
\[
 q=\frac{X_a+X_{a+1}}2,\qquad
 z=c(X_{a+1}-X_a)>0.
\]
Keep the single coordinate, if any, unchanged.
The coordinate Jacobian is \(c^{-M}\), so the corresponding
isometric change of \(L^2\) coordinates multiplies a wavefunction
by \(c^{-M/2}\).  Pair centers and the single coordinate are
called separated when all their mutual distances are positive.
The sets of nonseparated center coordinates have measure zero.
For \(\delta>0\) and \(Z<\infty\), a retained pair chart means the
part of these coordinates on which all center distances exceed
\(\delta\) and \(0<z_a<Z\) for every pair.
The corresponding regions for distinct pair assignments are
disjoint once \(c\) is sufficiently large in terms of \(\delta,Z\).
Convergence in these charts will always be followed by their
joint exhaustion as \(\delta\downarrow0\) and \(Z\uparrow\infty\).
It does not mean that the whole wavefunction has an \(L^2\) limit
under every possible pairing coordinate map.

\begin{lemma}[Compactness in pair coordinates]
\label{lem:bethe-pair-compactness}
Fix \(W_0>0\), \(n_p=2M+\ell\), and \(C<\infty\).
Let \(c_j\to\infty\) and let \(\psi_j\) be normalized vectors in
the form domain of the circle Hamiltonian such that
\[
 \left\langle\psi_j,
        (H_{c_j}+Mc_j^2/2)\psi_j\right\rangle\leq C.
\]
After passage to a subsequence, the localized functions in the
pair coordinates just defined converge in \(L^2\) on every fixed
retained chart.  For each seam and pair assignment these limits
are compatible under enlargement of the chart and have the form
\begin{equation}
 f_{\alpha,\mathfrak p}(q,v)\,
 \exp\left\{-\frac12\sum_{a=1}^M z_a\right\}
 \bigotimes_{a=1}^M
        (s_a^L\otimes s_a^R).
\label{eq:B12}
\end{equation}
Here \(v\) and its color are present only if \(\ell=1\); the
coefficient \(f_{\alpha,\mathfrak p}\) includes that color.
As \(\delta\downarrow0\) and \(Z\uparrow\infty\), the sum of the
squared norms of these chart limits tends to \(1\), counted with
the fixed square partition of unity.
There is no norm lost away from the pair coincidence manifolds
or at their intersections.  Inner products of two sequences
converging in this sense converge to the corresponding sums of
inner products of their limits.
\end{lemma}

\begin{proof}
The form localization identity gives
\[
 \sum_\alpha
 \left\langle\vartheta_\alpha\psi_j,
      (H_{c_j}+Mc_j^2/2)\vartheta_\alpha\psi_j\right\rangle
 =
 \left\langle\psi_j,
      (H_{c_j}+Mc_j^2/2)\psi_j\right\rangle
 +\sum_\alpha\|\,|\nabla\vartheta_\alpha|\psi_j\|^2.
\]
The error is bounded independently of \(j\), and every localized
term is nonnegative by Lemma~\ref{lem:bethe-line}.
It suffices, therefore, to prove the assertion for one localized
function on the line, with bounded norm and bounded shifted form
energy.  Fix a bounded chamber box containing the supports of
all these localized functions.  In what follows, restrict the
output of every reconstruction kernel to this box.  This is an
\(L^2\) contraction and changes none of the localized functions;
it is needed because a momentum-truncated reconstruction need
not itself have compact support.

Apply the isometric transform of Lemma~\ref{lem:bethe-line}.
Restrict all elementary momenta and pair centers to a fixed
bounded set, say absolute value at most \(R\).
The form identity bounds the squared norm discarded by this
restriction by \(C'/R^2\).  Its reconstruction has at most this
norm because \(\mathcal F_c^*\) is a contraction.
Likewise, all sectors with fewer than \(M\) pairs have
\(\lambda_{\ell'm'}\geq c^2/2\); their total squared transform
norm is at most \(2C'/c^2\).  Thus only the sector with \(M\)
pairs and \(\ell\) singles can survive, and its momenta may first
be kept in a fixed bounded set.

For that sector the reconstruction kernel contains the factor
\((2c)^{M/2}/\sqrt{\ell!\,M!}\).  Consider a group with overlapping
pair spans or with a single inside a pair span.
If a component contains \(p\) pairs and \(s_0\) singles, it has
\(2p+s_0\) slots and \(2p+s_0-1\) internal gaps.
For a component consisting of one adjacent pair this is one gap.
For an overlapping component or one containing an inserted
single it is at least \(p+1\) gaps.
Squaring \eqref{eq:B11} and integrating the ordered coordinates
therefore gives one factor \(O(c^{-1})\) for every such gap.
Over all components there are at least \(M+1\) decaying gaps in
these groups.  The squared kernel normalization is \(O(c^M)\),
so their squared Hilbert--Schmidt norms on the momentum cutoff
are \(O(c^{-1})\).  The remaining component positions range over
a fixed bounded set.  The bounded momentum volume and the
Cauchy--Schwarz inequality now show that all these groups have
vanishing reconstructed \(L^2\) norm, uniformly for bounded
transform coefficients.

Every surviving shuffle is a permutation of whole adjacent
blocks.  In coordinates for its adjacent-pair assignment, its
normal plane factor is
exactly \(\exp(-\sum z_a/2)\) and its color factor is the product
of double singlets in \eqref{eq:B12}.  The coefficients between
whole blocks are the explicit scalars
\eqref{eq:bethe-pair-single}--\eqref{eq:bethe-pair-pair}, whose
limits on a fixed momentum cutoff exist as \(c\to\infty\).
The rescaled Jacobian cancels the factor \(c^{M/2}\) in the
kernel normalization.  The fixed center-coordinate volume,
the finite momentum cutoff, and the exponential normal factor
give an integrable dominating square.
The reconstruction kernels for this assignment consequently
converge in Hilbert--Schmidt norm on every retained chart.
On a chart belonging to another assignment, their contribution
tends to zero: bounded rescaled gaps for both assignments would
force two distinct block centers to approach each other, whereas
the retained chart separates those centers by \(\delta\).
The exponential gap bound controls the remaining contribution.

This kernel convergence implies compactness even though the
transform coefficients depend on \(j\).  On the fixed cutoff
their \(L^2\) norms are bounded; extract a weakly convergent
subsequence.  The limiting Hilbert--Schmidt operator is compact,
and convergence to it in operator norm then gives strong
convergence of the reconstructed functions.
The limit has the asserted normal and color factors.

We verify that this argument retains the whole norm.
Restrict first to centers separated by at least \(\delta>0\)
and to \(0<z_a<Z\).  Different adjacent-pair assignments then
occupy disjoint configuration regions for all sufficiently
large \(c\).  The retained region is the union over those
assignments; each whole-block kernel is estimated in the
coordinates of its own assignment.
For the bounded-momentum kernels, the squared
norm excluded by center distances below \(\delta\) is uniformly
small as \(\delta\downarrow0\): the kernels are uniformly
bounded in the center coordinates, and these collision
neighborhoods have vanishing volume.
The norm excluded by \(z_a>Z\) tends uniformly to zero as
\(Z\to\infty\), by the exponential normal bound.
The groups already discarded cover the overlaps and inserted
singles, and have vanishing norm.
Thus neither intersections of pair coincidence manifolds nor
their complement retain additional norm.

Finally remove the momentum cutoff using the uniform
\(C'/R^2\) bound, and take a diagonal subsequence over the
finitely many seam charts and pair assignments and a countable
exhaustion of the separation and gap cutoffs.
The exact identity \(\sum_\alpha\vartheta_\alpha^2=1\)
then gives the asserted total norm.
Applying the same cutoffs to two sequences, strong convergence
on the retained charts gives convergence of their inner
products; Cauchy--Schwarz controls all discarded parts.
\end{proof}

The needed ground-state test is now an overlap argument.  The center
wave need not be characterized by a separate effective Hamiltonian:
strict positivity of its limiting shape suffices.

\begin{corollary}[Ground-state test by positive overlap]
\label{cor:bethe-ground-test}
Fix a circle length \(W_0>0\) and a particle number
\(n_p=2M+\ell\), with \(\ell\in\{0,1\}\), and work in the
fixed-weight sector balanced in each binary index.  Suppose a normalized eigenbranch
has energy \(-Mc^2/2+O(1)\) as \(c\to\infty\).  Suppose also that,
after one common phase choice for all charts, every subsequential limit
in the pair charts has the form \eqref{eq:B12}.  Require its scalar
center coefficient to be strictly positive almost everywhere that the
limiting chart cutoff is positive, in the static-sign convention of
Lemma~\ref{lem:brownian-ground}; when \(\ell=1\), the remaining single
has the majority color in each index.
Then the branch is the ground state for all sufficiently large \(c\).
\end{corollary}

\begin{proof}
Trial states consisting of separated, exponentially bound
pairs and smooth center wavefunctions give a circle ground-state
energy at most \(-Mc^2/2+O(1)\).  One can cut off the pair gaps
inside fixed disjoint intervals: the cutoff error is
exponentially small in \(c\), and the center kinetic energy is
bounded.  The true normalized ground state therefore satisfies
Lemma~\ref{lem:bethe-pair-compactness}.
It is the simple ground state in this sector, is nonnegative
in the static-sign convention, and its limits
retain total norm \(1\).  Their normal factors are
\eqref{eq:B12}; in a sector balanced in each binary index, the
remaining single, if present, has the majority color in each
index.  A nonzero nonnegative limit consequently has strictly positive
inner product with the strictly positive limiting center coefficient of
the eigenbranch.

If the two eigenstates were different for an unbounded sequence
of couplings, self-adjointness and the simplicity of the ground
state would make their inner products zero on that sequence.
Pass to a simultaneous compactness subsequence.  The inner-product
convergence just proved gives a strictly positive limit, a
contradiction.  This identifies the strong-coupling branch from
its wavefunction and energy, without assuming a string
description of the ground state on a finite circle.
\end{proof}
